% element: 10-s047:e04
% element: 10-s047:e05
\BeleskeID{10-s047}
Nagib pravolinijskog prelaza dobija se iz krajnjih tačaka karakteristike:
\[a=\frac{V_{CC}-V_{CES}}{(V_D+V_{\gamma T})-(2V_D+V_{BES}-V_{\gamma D})}.\]
Zamena tog odnosa u jednačinu prave daje
\[V_O=\frac{V_{CC}-V_{CES}}{(V_D+V_{\gamma T})-(2V_D+V_{BES}-V_{\gamma D})}\bigl(V_I-(V_D+V_{\gamma T})\bigr)+V_{CC}.\]
Prag preslikavanja samog u sebe određujemo iz $V_I=V_O$:
\[V_I=a\bigl(V_I-(V_D+V_{\gamma T})\bigr)+V_{CC}.\]
Rešavanjem po $V_I$ dobija se $V_M$.

Margine za višestruke izvore su $V_{IL}-V_{OL}$ i $V_{OH}-V_{IH}$, a za jedan izvor $V_M-V_{OL}$ i $V_{OH}-V_M$. Zbog uske prelazne zone njihove vrednosti biće slične.
